% Source: physical PDF pp. 211--215 (printed pp. 188--192), from the Section
% 28.2 heading through Eq. (28.2.19), immediately before Section 28.3.
% Inventory: Eqs. (28.2.1)--(28.2.19), no unnumbered display groups, linked
% citations [5], [6], [7], and [7a], Table 28.1, and no figures, footnotes,
% or centered three-asterisk dividers.
% Corrected a mistyped equation reference and ``affect''/``effect''.
% Uncertainties: none.

\subsection{Supersymmetry and Strong--Electroweak Unification}

We shall have to defer a detailed assessment of supersymmetric models of
particle physics until we are ready to consider how supersymmetry is broken.
In this section we will consider the quantitative application of
supersymmetry in one context in which the mechanism for the breakdown of
supersymmetry is relatively unimportant, and in which supersymmetry has
scored what so far is its greatest empirical success.

If the \(SU(3)\times SU(2)\times U(1)\) gauge group of the strong and
electroweak interactions is embedded in a simple group \(G\) that has the
known quarks and leptons (plus perhaps some
\(SU(3)\times SU(2)\times U(1)\)-neutral fermions) as a representation, then,
as described in Section 21.5, at energies at or above the scale \(M_X\) at
which \(G\) is spontaneously broken, the
\(SU(3)\times SU(2)\times U(1)\) coupling constants will be related by
\begin{equation}
  g_s^2=g^2=\frac{5g'^2}{3}
  \qquad\text{at energies }\geq M_X\,.
  \tag{28.2.1}\label{eq:28.2.1}
\end{equation}
At energies far below \(M_X\), these couplings are seriously affected by
renormalization corrections. If measured at a scale \(\mu<M_X\), the
couplings will have values \(g'^2(\mu)\), \(g^2(\mu)\), \(g_s^2(\mu)\),
governed by one-loop renormalization group equations
\begin{equation}
  \mu\frac{\dd}{\dd\mu}g'(\mu)=\beta_1\bigl(g'(\mu)\bigr)\,,
  \qquad
  \mu\frac{\dd}{\dd\mu}g(\mu)=\beta_2\bigl(g(\mu)\bigr)\,,
  \qquad
  \mu\frac{\dd}{\dd\mu}g_s(\mu)=\beta_3\bigl(g_s(\mu)\bigr)\,,
  \tag{28.2.2}\label{eq:28.2.2}
\end{equation}
with initial conditions at \(M_X\) satisfying Eq.~\eqref{eq:28.2.1}. In the
original use~\hyperref[ch28-ref-5]{[5]} of these renormalization group
equations, discussed in Section 21.5, the beta functions were calculated in
one-loop order to be
\begin{equation}
  \beta_1=\frac{5n_g g'^3}{36\pi^2}\,,
  \tag{28.2.3}\label{eq:28.2.3}
\end{equation}
\begin{equation}
  \beta_2=\frac{g^3}{4\pi^2}
  \left(-\frac{11}{6}+\frac{n_g}{3}\right)\,,
  \tag{28.2.4}\label{eq:28.2.4}
\end{equation}
\begin{equation}
  \beta_3=\frac{g_s^3}{4\pi^2}
  \left(-\frac{11}{4}+\frac{n_g}{3}\right)\,,
  \tag{28.2.5}\label{eq:28.2.5}
\end{equation}
where \(n_g\) is the number of generations of quarks and leptons and the
relatively small contributions of scalar fields are here neglected. Since
\(M_X\) will turn out to be many orders of magnitude larger than the
energies accessible with today's accelerators, it seems reasonable to
suppose that supersymmetry is unbroken over most of the range below \(M_X\),
in which case all of the new fields discussed in the previous section need
to be included in calculations of the beta functions in
Eq.~\eqref{eq:28.2.2}. These new fields introduce three major changes in the
calculations of the beta functions:

\noindent\textbf{1.}\quad
For every gauge boson, there is a Majorana gaugino with the same
\(SU(3)\times SU(2)\times U(1)\) quantum numbers. Equation
\eqref{eq:17.5.41} shows that the ratio of the contribution to the beta
function for any gauge coupling of a Dirac fermion that furnishes a
representation of the gauge group with generators \(t_A\) to the
contribution of the corresponding gauge boson is \(-4C_2/11C_1\), where
according to Eqs.~\eqref{eq:17.5.33} and \eqref{eq:17.5.34} the ratio of
\(C_1\) and \(C_2\) is given by:
\begin{equation}
  \sum_{AB}C_{CAB}C_{DBA}
  =-(C_1/C_2)\operatorname{Tr}(t_Ct_D)\,.
  \tag{28.2.6}\label{eq:28.2.6}
\end{equation}
For the adjoint representation, \((t_C)_{AB}=\ii C_{ABC}\), so \(C_1=C_2\),
and so a Dirac fermion in the adjoint representation makes a contribution
that is \(-4/11\) that of the gauge bosons. But the gauginos are Majorana
fermions, so their contribution is \(-2/11\) that of the gauge boson. Thus
the term \(11/6\) and \(11/4\) in Eqs.~\eqref{eq:28.2.4} and
\eqref{eq:28.2.5} are reduced by a factor \(9/11\) to \(9/6\) and \(9/4\),
respectively.

\noindent\textbf{2.}\quad
For every left-handed quark, lepton, antiquark, or antilepton field, there
is a complex scalar field with the same
\(SU(3)\times SU(2)\times U(1)\) quantum numbers. Following the same method
as in Section 17.5, it is not hard to calculate that the contribution of a
complex scalar field belonging to a representation of a gauge group with
generators \(t_A\) to the beta function for the gauge coupling \(g_i\) is
\begin{equation}
  \bigl[\beta_i(g_i)\bigr]_{\mathrm{scalar}}
  =\frac{g_i^3C_{2\ii}}{48\pi^2}\,,
  \tag{28.2.7}\label{eq:28.2.7}
\end{equation}
where \(\operatorname{Tr}(t_At_B)=g_i^2C_{2\ii}\delta_{AB}\). This is \(1/4\)
the contribution of a Dirac spinor field in the same representation, given
by Eq.~\eqref{eq:18.7.2}, and hence \(1/2\) the contribution of each
left-handed spinor field (including the complex conjugates of the
right-handed components of the Dirac fields). Thus the coefficients of
\(n_g\) in Eqs.~\eqref{eq:28.2.3}--\eqref{eq:28.2.5} should be increased
by a factor \(3/2\).

\noindent\textbf{3.}\quad
The decrease by a factor \(9/11\) of the negative gauge boson terms in the
beta function and the increase by a factor \(3/2\) of the positive squark
and slepton terms both lead to a general decrease in the rate at which the
three gauge coupling constants diverge below \(M_X\) from the ratios
\eqref{eq:28.2.1}. This will increase our estimate of \(M_X\), but, as we
shall see, in itself it would have no effect on the prediction of the
electroweak mixing parameter \(\sin^2\theta\). But these changes do enhance
the relative contribution of the Higgs scalars, which was neglected in
Eqs.~\eqref{eq:28.2.3}--\eqref{eq:28.2.5}, and which is now also
accompanied with the larger contribution of the accompanying higgsinos.
With \(n_s\) of the superfields \((H_1^0,H_1^-)\) or
\((H_2^+,H_2^0)\) discussed in the previous section, the constant
\(C_{2\ii}\) in Eq.~\eqref{eq:28.2.7} is
\([(1/2)^2+(-1/2)^2]n_s=n_s/2\) for \(SU(2)\), and is
\(2n_s(\pm1/2)^2=n_s/2\) for \(U(1)\). According to
Eq.~\eqref{eq:28.2.7}, the scalar components of these superfields make a
contribution to \(\beta_1\) equal to \(n_sg'^3/96\pi^2\) and a contribution
to \(\beta_2\) also equal to \(n_sg^3/96\pi^2\). As we have seen, the
Majorana higgsinos contribute twice as much to the beta functions as
complex scalars with the same quantum numbers, so the superfields
\((H_1^0,H_1^-)\) or \((H_2^+,H_2^0)\) make a total contribution to
\(\beta_1\) and \(\beta_2\) that is \(3/2\) the contribution of the Higgs
scalars, and hence equal to \(n_sg'^3/32\pi^2\) and \(n_sg^3/32\pi^2\),
respectively.

Making all these changes in the beta functions, we now have
\begin{equation}
  \beta_1=\frac{g'^3}{4\pi^2}
  \left(\frac{5n_g}{6}+\frac{n_s}{8}\right)\,,
  \tag{28.2.8}\label{eq:28.2.8}
\end{equation}
\begin{equation}
  \beta_2=\frac{g^3}{4\pi^2}
  \left(-\frac{9}{6}+\frac{n_g}{2}+\frac{n_s}{8}\right)\,,
  \tag{28.2.9}\label{eq:28.2.9}
\end{equation}
\begin{equation}
  \beta_3=\frac{g_s^3}{4\pi^2}
  \left(-\frac{9}{4}+\frac{n_g}{2}\right)\,.
  \tag{28.2.10}\label{eq:28.2.10}
\end{equation}
The solutions of the renormalization group equations \eqref{eq:28.2.2} are
then
\begin{equation}
  \frac{1}{g'^2(\mu)}
  =\frac{1}{g'^2(M_X)}
  +\frac{1}{2\pi^2}
  \left(\frac{5n_g}{6}+\frac{n_s}{8}\right)
  \ln\left(\frac{M_X}{\mu}\right)\,,
  \tag{28.2.11}\label{eq:28.2.11}
\end{equation}
\begin{equation}
  \frac{1}{g^2(\mu)}
  =\frac{1}{g^2(M_X)}
  +\frac{1}{2\pi^2}
  \left(-\frac{3}{2}+\frac{n_g}{2}+\frac{n_s}{8}\right)
  \ln\left(\frac{M_X}{\mu}\right)\,,
  \tag{28.2.12}\label{eq:28.2.12}
\end{equation}
\begin{equation}
  \frac{1}{g_s^2(\mu)}
  =\frac{1}{g_s^2(M_X)}
  +\frac{1}{2\pi^2}
  \left(-\frac{9}{4}+\frac{n_g}{2}\right)
  \ln\left(\frac{M_X}{\mu}\right)\,.
  \tag{28.2.13}\label{eq:28.2.13}
\end{equation}
It is convenient to take \(\mu=m_Z\), so that \(SU(2)\times U(1)\) can be
regarded as unbroken over almost all of the range of energies in which we
use the formulas \eqref{eq:28.2.11}--\eqref{eq:28.2.13}. Using
Eq.~\eqref{eq:28.2.1}, the difference between Eqs.~\eqref{eq:28.2.12} and
\eqref{eq:28.2.13} gives
\begin{equation}
  \frac{1}{g^2(m_Z)}-\frac{1}{g_s^2(m_Z)}
  =\frac{1}{2\pi^2}
  \left(\frac{3}{4}+\frac{n_s}{8}\right)
  \ln\left(\frac{M_X}{m_Z}\right)\,,
  \tag{28.2.14}\label{eq:28.2.14}
\end{equation}

\begin{table}[H]
  \renewcommand{\thetable}{28.1}
  \centering
  \caption{Values of the electroweak mixing parameter \(\sin^2\theta\) and
  unification mass \(M_X\) given by Eqs.~\eqref{eq:28.2.17} and
  \eqref{eq:28.2.18}, as functions of the number \(n_s\) of left-chiral
  superfield doublets \((H_1^0,H_1^-)\) or \((H_2^+,H_2^0)\).}
  \label{tab:28.1}
  \begin{tabular}{c@{\qquad}c@{\qquad}c}
    \hline\hline
    \(n_s\) & \(\sin^2\theta\) & \(M_X\) (GeV) \\
    \hline
    \(0\) & \(0.203\) & \(8.7\cdot10^{17}\) \\
    \(2\) & \(0.231\) & \(2.2\cdot10^{16}\) \\
    \(4\) & \(0.253\) & \(1.1\cdot10^{15}\) \\
    \hline\hline
  \end{tabular}
\end{table}

while the difference between Eq.~\eqref{eq:28.2.12} and \(3/5\) of
Eq.~\eqref{eq:28.2.11} gives
\begin{equation}
  \frac{1}{g^2(m_Z)}-\frac{3}{5g'^2(m_Z)}
  =\frac{1}{2\pi^2}
  \left(-\frac{3}{2}+\frac{n_s}{20}\right)
  \ln\left(\frac{M_X}{m_Z}\right)\,.
  \tag{28.2.15}\label{eq:28.2.15}
\end{equation}
Equation \eqref{eq:21.3.19} allows us to express the electroweak couplings in
terms of the electroweak mixing angle \(\theta\) and the positron charge
\(e\):
\begin{equation}
  g(m_Z)=-e(m_Z)/\sin\theta\,,
  \qquad
  g'(m_Z)=-e(m_Z)/\cos\theta\,.
  \tag{28.2.16}\label{eq:28.2.16}
\end{equation}
We can then solve for the unknowns \(\ln(M_X/m_Z)\) and \(\sin^2\theta\) in
terms of input parameters \(e(m_Z)\) and \(g_s(m_Z)\):
\begin{equation}
  \sin^2\theta
  =
  \frac{
    18+3n_s+\bigl(e^2(m_Z)/g_s^2(m_Z)\bigr)(60-2n_s)
  }{108+6n_s}\,,
  \tag{28.2.17}\label{eq:28.2.17}
\end{equation}
\begin{equation}
  \ln\left(\frac{M_X}{m_Z}\right)
  =
  \left(\frac{8\pi^2}{e^2(m_Z)}\right)
  \left(
    \frac{1-\bigl(8e^2(m_Z)/3g_s^2(m_Z)\bigr)}{18+n_s}
  \right)\,.
  \tag{28.2.18}\label{eq:28.2.18}
\end{equation}
For \(n_s=0\) Eq.~\eqref{eq:28.2.17} gives the same result
\eqref{eq:21.5.15} for \(\sin^2\theta\) as was originally calculated
(ignoring the small contribution of Higgs scalars) in non-supersymmetric
theories, but the value \eqref{eq:28.2.18} of \(\ln(M_X/m_Z)\) is larger
than the original result \eqref{eq:21.5.16} by a factor \(11/9\), which as
we have seen arises from gaugino contributions to the beta functions.

Using the same input parameters \(e^2(m_Z)/4\pi=(128)^{-1}\),
\(g_s^2(m_Z)/4\pi=0.118\), \(m_Z=91.19\ \mathrm{GeV}\) as in Section 21.5
now gives the numerical results shown in Table~\ref{tab:28.1}. As discussed
in the previous section, the necessity of cancelling anomalies in the
electroweak currents requires equal numbers of \((H_1^0,H_1^-)\) and
\((H_2^+,H_2^0)\) doublets, so we consider only even values of the number
\(n_s\) of these superfields.

Remarkably, the value \(n_s=2\) for the simplest plausible theory yields a
value~\hyperref[ch28-ref-6]{[6]} \(\sin^2\theta=0.231\) which is in perfect
agreement with the experimentally observed value, \(\sin^2\theta=0.23\).
The value of \(M_X\) is 20 times
greater~\hyperref[ch28-ref-7]{[7]} than calculated in this way in
non-supersymmetric theories, leading to a decrease by a factor \(20^{-4}\)
in the rate for proton decay processes like
\(p\longrightarrow\pi^0+e^+\), thus removing a conflict with the
experimental non-observation of such processes. (Proton decay is discussed
in more detail in Section 28.7.) This increase in the value of \(M_X\)
brings it closer to the energy scale \(\approx10^{18}\ \mathrm{GeV}\) at
which gravitation has the same strength as other interactions. It may be
that this remaining gap may be filled by a change in gravitational
interactions at very high energies~\hyperref[ch28-ref-7a]{[7a]}.

A value \(n_s=4\) would give a value for \(\sin^2\theta\) in serious
disagreement with experiment, and a value of \(M_X\) low enough to revive
the conflict with expectations for proton decay. This makes a strong case
for having just one of each superfield \((H_1^0,H_1^-)\) and
\((H_2^+,H_2^0)\).

Unlike the calculated values of \(\sin^2\theta\) and \(M_X\), the calculated
value of the common gauge coupling \eqref{eq:28.2.1} at \(M_X\) does depend
on the number of generations as well as the number of scalar doublets. With
\(n_g=3\) and \(n_s=2\) and our previous input parameters,
Eq.~\eqref{eq:28.2.13} gives
\begin{equation}
  \frac{g^2(M_X)}{4\pi}
  =\frac{g_s^2(M_X)}{4\pi}
  =\frac{1}{17.5}\,.
  \tag{28.2.19}\label{eq:28.2.19}
\end{equation}
